% ============================================================
%  同济大学 2027 届 本科毕业设计（论文）开题报告
%  课题：面向未来混合自动驾驶环境的风险传播机理研究
%  学生：禹尧珅（2352396）  指导教师：李王明卉
%
%  ── 怎么编 ──────────────────────────────────────────────
%  编译器选 XeLaTeX（不要 pdfLaTeX，排不了中文）。
%  命令行：latexmk -xelatex main.tex   或   xelatex main.tex 连跑两遍
%  编辑器里把编译器设成 XeLaTeX 就能直接编，不需要配任何路径。
%
%  ── 下面这一段别删 ──────────────────────────────────────
%  文档类 tongjithesis.cls 在 style/ 子目录，字体定义在 style/font/。
%  TeX 默认不去子目录找文件，用 \input@path 告诉它去哪儿找，
%  这样任何编辑器、任何工作目录都能编过。
% ============================================================
\makeatletter
\def\input@path{{style/}{style/font/}}
\makeatother

\documentclass[oneside,doctype=proposal,fontset=windows,times=true]{tongjithesis}

% ============================================================
%  封面信息（要改就改这里）
%  全部取自已定稿的《2027 届本科毕业论文任务书》，不另编造；
%  与任务书不一致时以任务书原件为准。
% ============================================================
\school{电子与信息工程学院}
\major{人工智能}
\student{2352396}{禹尧珅}
\thesistitle{面向未来混合自动驾驶环境的风险传播机理研究}{}
\thesistitleeng{Research on Risk Propagation Mechanisms in Future Mixed Autonomous Driving Environments}{}
\thesisadvisor{李王明卉}
\thesisdate{\the\year}{\the\month}{\the\day}
\taskbookperiod{2026}{9}{14}{2027}{4}{28}

% ── 参考文献：BibTeX + biber（文档类已载入 biblatex style=gb7714-2025）──
%    正文写 \cite{key}，编号由 biber 现算，不手写死；条目在 refs.bib。
\hypersetup{colorlinks=true,linkcolor=black,citecolor=black,urlcolor=blue}
\usepackage{tikz}
\usetikzlibrary{positioning,arrows.meta}
\usepackage{xurl}   % DOI 与 URL 允许在任意位置断行，避免数字被拆开
\addbibresource{refs.bib}

% ── 标点悬挂：与 Word「允许标点溢出边界」一致，行末标点可探出版心 0.5 字 ──
\XeTeXprotrudechars=2

\begin{document}
\pagestyle{empty}

% 封面
\MakeDocument

% ============================================================
%  正文：与官方《2毕业设计(论文)开题报告.doc》逐段对应（2026-11-27 按 Word COM
%  逐段读出的原始属性重建，段落顺序、段前段后、行距、缩进全部照抄，不靠挪坐标）。
%
%  ⚠ 现在正文是「空白模板」状态——和官方 Word 原件一模一样：小节下面就是一个空段落。
%    编译出来的 PDF 应当与官方 PDF 逐页、逐行重合（3 页）。
%
%  ✍ 怎么写：把标注「▼ 写在这里」的空段落 \tjformblank 换成你的正文段落
%    （用 tjformbody 环境，宋体小四、首行缩进 2 字符、1.5 倍行距）。
%    每个空段落下方都留了「这一段该写什么」的提示，提示只是注释，不会印出来。
%
%  🔴 防泄露（导师 2026-11 微信要求）：正文不得写 AAD 论文未公开的核心创新点——
%     消息可达性判据的具体定义、H1–H4 因果识别框架、十四类消息生命周期分类、
%     区域风险聚合的具体公式、现有量化结论。只写研究问题、技术路线、公开指标。
% ============================================================

\tjformsection{一、\tjheadertext 课题背景}
\tjformgap{18.72bp}%
\tjformsubsection{1．选题背景与研究意义}
\begin{tjformbody}
自动驾驶车辆的渗透率难以在短期内快速提升。依据当前的技术演进与分级推进节奏，未来较长时期内，同一路段将同时运行控制方式与网联能力各不相同的车辆。本文按控制方式与网联能力，将研究对象区分为人工驾驶车辆、配备辅助驾驶系统的人工驾驶车辆以及具备车车通信能力的高等级自动驾驶车辆。受制于存量车辆的更新周期，混合交通将长期延续。三者的差异主要源于对同一交通扰动的响应机制，车辆本体的动力学性能并非主要决定因素。人工驾驶依赖驾驶员的视觉感知与经验判断；辅助驾驶系统的感知范围受限于车载传感器视距；网联自动驾驶车辆则可借助车车通信获取超视距信息。三类车辆混编行驶时，交通流运行规律不再均一，车流内部的交互复杂度随之上升。

高密度车流会迅速放大上述响应差异。三类车辆在感知范围、决策逻辑、响应时延及跟驰换道行为上存在系统性分化。人工驾驶员的制动反应伴随可观测的延迟，而在通信链路可靠且控制器有效接收信息的条件下，网联自动驾驶车辆可借助超视距信息辅助决策，其响应效果仍取决于通信时延与控制策略。一旦车流中某车辆实施紧急制动或强制换道，后方车辆的减速幅度与时机便会出现错位。前车异常动作首先触发紧邻后车的避险制动，继而沿车流向上游传导，最终在较远的上游路段诱发二次追尾或局部拥堵。风险因而不再局限于单车范围，沿车队向上游扩散。

既有安全评价方法难以刻画此类跨车辆扩散过程。现有研究多以单车或车辆对为基本评价单元，借助碰撞时间（TTC）、后侵入时间（PET）等替代安全指标度量危险程度\cite{minderhoud2001}\cite{tarko2018}\cite{arun2021}\cite{wang2021}。此类指标可较准确地判别相邻车辆对在特定时刻的危险状态，却不足以表征风险沿车流的传递过程。其隐含假设将相邻车辆对的交互视为独立事件，不描述车辆间的连续响应，难以表征下游扰动向上游传播时对车群跟驰间距与速度分布的影响。当车流中同时存在网联消息广播与物理制动波时，基于车辆对的指标既无法刻画多车连续响应所累积的区域性风险，亦难以评估远端预警信息对下游车群的干预成效。

评价视角的局限暴露出两方面研究不足。其一为量化框架的欠缺：远端风险信息、车辆动作调整与区域总体风险三者之间，尚缺乏统一的量化分析框架。其二为两类到达路径的混同：后车减速既可能源于对前车制动灯及间距压缩的视觉响应，亦可能源于对远端车辆预警消息的接收，既有研究在剖析风险传播时缺乏区分物理影响与信息影响的手段。两类路径若不加区分，车辆动作改变的归因便难以确立，网联通信对风险传播的抑制或放大效应亦无从定量评估。

填补上述缺口对混合交通安全研究具有直接意义。目前，物理冲击与网联信息在异质车流中的传递规律尚缺乏系统性认知，厘清该规律有助于阐明风险在何种渗透率下被放大、在何种条件下趋于衰减，从而贯通微观驾驶行为与宏观车流安全之间的分析链条。突破车辆对评价的局限后，跨车流风险追踪将具备量化基础，路侧管控与车路协同预警的触发时机亦可获得依据，研究结果可为评估路侧管控与车路协同预警策略提供分析依据，其对实际事故风险的影响仍需进一步验证。

据此，本课题聚焦的核心科学问题为：在混合自动驾驶环境中，来源可追溯且确被控制器采纳的远端风险消息，能否在物理影响抵达目标区域之前，通过改变车辆动作而重塑区域总体风险。研究以微观交通仿真为主要手段，以消息接收状态作为干预变量，先检验其是否引起车辆动作变化，再评估该干预对区域总体风险的影响。
\end{tjformbody}

\tjformsubsection{2．国内外研究现状（文献综述）}
\begin{tjformbody}
交通安全评价长期依赖事故统计数据，然而事故样本稀疏且统计周期漫长，难以支撑混合交通流的风险推演。为此，学界引入替代安全指标，其中 TTC、减速度需求（DRAC）与 PET 的应用最为广泛。TTC 度量两车维持当前速度时的碰撞剩余时间，DRAC 计算后车避免追尾所需的最小减速度，PET 则记录前车驶离冲突区至后车进入该区域的时间差\cite{minderhoud2001}\cite{tarko2018}\cite{arun2021}。三类指标将离散的事故结果转化为可连续观测的运动学状态。

上述指标各自存在明确的适用边界。TTC 与 DRAC 主要针对同车道跟驰追尾场景；当后车速度低于前车时，TTC 趋于无穷大而丧失判别能力。PET 适用于交叉口或换道交织区的轨迹交叉冲突，却无法反映未发生轨迹重叠时的纵向逼近风险\cite{arun2021}\cite{wang2021}。在混合自动驾驶环境下，人工驾驶车辆与自动驾驶车辆在反应时间、制动特性及跟驰间距上存在明显差异，固定阈值对不同车辆类型与运行条件的适用性可能受限，其误报与漏报情况需通过实证或仿真分析评估\cite{wang2021}\cite{fu2021}。

替代指标的计算大多局限于单车或车辆对层面，仅能刻画局部瞬时冲突。当某车辆因前车紧急制动而采取避险动作时，该局部冲突会向后方车流传导\cite{li2014}。单车指标无法捕捉这一传导过程，亦难以反映网联信息介入后风险在车队中的放大或衰减效应，评价标度亟需由车辆对扩展至车队乃至路段。

风险在车队中的传递与交通扰动的传播规律密切相关。早期研究多采用宏观流体力学模型，将车流视为连续介质，借助激波理论解释拥堵波的上游传播\cite{lighthill1955}\cite{richards1956}\cite{newell1993}。此后，微观跟驰模型被引入车队弦稳定性分析，微观仿真中的车辆纵向运动亦普遍由该类模型驱动\cite{treiber2000}，相关研究分析了在特定模型与参数条件下，微小速度扰动沿人工驾驶车队逐车传播并可能诱发拥堵的过程。此类工作为扰动传播构建了动力学描述，但其研究对象主要局限于同质车流。

近年来，高精度轨迹数据与路侧感知设备的普及推动了数据驱动的风险传播研究。部分学者利用无人机或固定摄像头提取车辆轨迹，结合复杂网络或图模型识别路网中的风险热区与传播路径\cite{saberi2020}\cite{duan2023}\cite{hamedmoghadam2021}。此类方法能够从海量历史数据中挖掘高风险路段与时空关联特征，为传播过程的实证刻画提供了有效手段。

然而，宏观流体模型抹平了个体驾驶行为的异质性，难以刻画自动驾驶车辆与人工驾驶车辆混行时的交互差异。数据驱动模型虽可拟合风险热区，却往往缺乏行为机理层面的解释，既无法回答特定预警消息如何改变下游风险分布这类因果问题，亦难以向未观测到的高渗透率场景外推。本研究拟将微观跟驰动力学与时空传播分析相结合，旨在刻画混合车流中个体避险动作向区域风险演化的传导路径与衰减规律。

上述传导路径的验证有赖于微观交通仿真平台。SUMO、VISSIM 等仿真工具已被广泛用于构建混合交通测试场景，其中 SUMO 凭借开源特性与良好可扩展性获得大量研究采用\cite{sumo2010}。研究者通常基于实测轨迹数据集对跟驰模型与换道模型进行参数标定，NGSIM、highD 与 pNEUMA 是较常采用的公开数据集，其道路类型、采样频率与轨迹完整性各有差别\cite{coifman2017}\cite{barmpounakis2020}；通过调整最大减速度、期望车头时距与反应时间等参数，使仿真车流的宏观基本图与微观车头时距分布贴近真实道路状态\cite{cunto2008}。

车路协同信息接入混合交通流后，仿真建模的复杂度随之上升。通信约束重塑了车辆的感知边界，已有研究在仿真器中耦合通信网络模块，或直接引入通信时延与丢包率参数，以模拟网联车辆获取前车状态或路侧预警的过程\cite{orosz2016}。部分工作亦探讨了不同自动驾驶渗透率下，网联信息对车队通行能力与冲突频次的影响。

在网联信息建模环节，现有仿真研究多将其简化为理想广播或固定时延，忽略了消息在复杂交通流中的实际接收状态与时效衰减。同时，仿真中较少区分人工驾驶员与自动驾驶系统在接收消息后的差异化响应逻辑。本研究拟构建将通信时延、丢包与车辆异质性响应耦合于同一闭环的仿真环境，并利用实测轨迹数据完成基础交通流参数的标定与校核。

通过主动干预抑制风险传播构成另一条研究脉络。现有风险抑制手段主要包括路侧可变限速、动态路径诱导、车内碰撞预警以及协同自适应巡航（CACC）等\cite{hegyi2010}\cite{khondaker2015}\cite{ploeg2011}\cite{milanes2014}。预警与诱导侧重于改变驾驶员的认知与路径选择，可变限速与 CACC 则直接调控车辆的纵向运动状态与车队间距。

评价上述干预措施时，学界普遍关注安全性、通行效率与交通流平稳性三类指标的权衡。例如，过早或过于频繁的碰撞预警虽可降低追尾风险，却可能诱发驾驶员过度制动，致使后方车流速度剧烈波动并削弱路段通行能力\cite{khondaker2015}。在混合交通流中，CACC 车队若因人工驾驶车辆频繁切入而退化为普通跟驰，其平稳性收益亦会明显下降。

从本文梳理的文献看，控制与消息机制方向的工作较多关注局部车队的弦稳定性优化或宏观路段的流量均衡\cite{stern2018}\cite{coifman2026}，对单条干预消息经由车辆动作影响区域风险的完整过程，仍有进一步分析的空间。消息设计不当甚至会诱发二次风险传播，例如局部急刹预警可能引发下游车道的连锁换道冲突。本研究拟在混合交通仿真环境中对比不同消息触发与控制策略下安全性、效率与平稳性的演变，识别能够有效阻断风险传播的干预条件与消息触发边界。

上述四类工作在单车冲突判别、宏观扰动传播、仿真平台构建以及局部控制优化方面已形成相当积累，但各环节之间仍存在明显断层。现有研究多停留于单车层面的替代安全指标计算，或直接面向宏观路段的流量与事故统计，从消息到动作、再到区域风险的中间传导环节缺乏可测量的刻画。网联消息如何改变个体车辆的决策动作，个体动作的叠加又如何重塑区域层面的风险分布，其机理尚不明晰。本研究拟在混合自动驾驶仿真环境中贯通追踪消息触发、车辆响应与风险传播三个阶段，检验远端消息是否被车辆控制器实际采纳，并评估其对区域总体风险的作用。
\end{tjformbody}

\tjformgap{31.47bp}%
\tjformsection[1.5]{二、\tjheadertext 方案介绍}
\tjformsubsection{1．主要研究内容}
\begin{tjformbody}
混合交通流中，人工驾驶车辆、辅助驾驶车辆与网联自动驾驶车辆将长期共存。既有研究多聚焦于宏观通行能力或单车碰撞概率，对微观风险在车队内部的传递路径尚缺乏精细刻画\cite{wang2021}。本研究以局部扰动向车队上游及相邻车道的扩散范围作为分析对象，不预设其必然扩散或必然衰减。

边界明确之后，首要任务在于量化单车风险并追踪其跨车传递。传统碰撞次数统计在短时窗口内样本稀疏，难以捕捉未发生碰撞的险态演化。本研究选用减速度需求（DRAC）作为核心度量指标，并结合实测制动分布校准风险阈值\cite{tarko2018}\cite{fu2021}。在此基础上，建立风险事件在时间与空间维度上的连续性识别方法，同时针对时空容差参数开展敏感性检验，以剔除偶发制动噪声。

真实道路中的高风险冲突样本稀少，仅凭实测数据难以覆盖多变的混合交通工况。仿真环境依托高分辨率真实轨迹数据标定跟驰与换道模型参数\cite{cunto2008}，并在微观交通仿真软件中构建测试场景。场景构建以单车道跟驰为起点，逐步扩展至包含强制换道与汇入行为的多车道情形，通过调整自动驾驶渗透率，复现不同车流结构下的典型冲突源事件。

场景就绪后，研究重心转向风险波在混合车流中的演化规律。实验通过系统调整自动驾驶渗透率、期望车头时距及网联通信延迟等参数，定量测度风险波的传播速度与衰减率，并对比不同工况下的风险扩散范围，进而识别诱发风险放大的关键参数组合。

针对高风险传播链，抑制策略从车辆控制与交通管理两个维度切入。控制侧重点测试网联车辆采纳远端风险预警消息后的制动平滑效果，评估其对上游风险波的阻断能力\cite{coifman2026}；管理侧依据仿真实验结果设计车速引导或车道管控策略，并通过对照实验评估其安全性、通行效率与运行平稳性\cite{khondaker2015}。
\end{tjformbody}

\tjformsubsection{2．研究方案}
\begin{tjformbody}
实证与标定的依据取自公开的高分辨率车辆轨迹数据集。所选数据具备较高的采样频率与完整的车辆运动学记录，NGSIM 的高采样率与 pNEUMA 的大范围覆盖各具优势\cite{coifman2017}\cite{barmpounakis2020}。数据处理遵循由简入繁的原则：先提取单车道无换道片段，获取位置、速度、加速度及车头间距等变量，用于标定跟驰参数与基线运动学特征；待单车道规律明确，再提取包含换道切入与汇入行为的多车道片段，补充横向交互特征。全部轨迹数据均经平滑滤波与异常值剔除，以确保速度与加速度导数的计算精度\cite{coifman2024}。

清洗后的轨迹数据用于构建以 DRAC 为核心的替代安全指标体系。该指标表征在前车速度保持不变、后车立即采取恒定减速度制动的假设下，后车避免追尾所需的最小减速度。研究利用前后车的瞬时速度差与净间距逐帧计算指标值，统计实测数据中的紧急制动分布，预先定义风险阈值的校准准则，并检验不同车辆类别下阈值设定的敏感性\cite{fu2021}。为区分孤立急刹与真实的风险传播，建立风险事件在时间与空间上的连续性识别方法，并针对时空容差参数开展敏感性检验，以验证识别结果的稳健性。

指标体系确立后，搭建自研模型与微观仿真软件耦合的测试平台。人工驾驶与自动驾驶的跟驰、换道逻辑由外部脚本独立编写，经由外部步进接口接入 SUMO，以 0.1 秒为步长接管车辆的逐帧运动控制\cite{sumo2010}。鉴于接口通信或积分步长可能引入数值偏差，平台强制执行同源数据回放核对：将相同的初始状态与前车轨迹同时输入独立脚本与耦合平台，研究预先设定轨迹位置与速度的误差指标及其容限，并据此检验两者的回放结果。

平台通过核对后，采用配对对照实验设计以剥离随机干扰。各实验条件采用配对随机种子并重复运行，确保背景车流的生成序列与初始扰动一致，报告重复次数与关键指标的统计量。实验设置多臂对照，分别比较纯人工车流、无网联混合车流与网联混合车流，并成组调整自动驾驶渗透率与源事件制动强度，观测风险波的响应差异。为防止观测窗口截断引发指标失真，另在不同时间窗口长度下开展跨窗口复核。

在多臂框架下进一步引入网联控制干预。控制侧实验在其余条件相同的前提下，按配对种子分配网联车辆是否接收远端风险预警消息，每条消息均关联可追溯的源头事件。以物理影响抵达目标区域之前的时段作为评价窗口，比较接收组与对照组的车辆动作与区域风险，估计该窗口内消息干预的效应，并将车辆动作变化作为解释该效应的机制证据。

上述环节共同构成完整的技术路线，如图 1 所示。研究以高分辨率真实轨迹数据为起点，完成数据清洗与运动学特征提取，进而构建并校准以 DRAC 为核心的指标体系与风险事件识别方法；随后依托经同源回放核对的耦合平台，开展配对种子与多臂对照实验，刻画风险传播规律；在此基础上引入可追溯的远端风险消息开展控制侧干预测试，最终归纳混合交通流的风险传播机理与管控启示。

\begin{center}
\begin{tikzpicture}[
  font=\zihao{5}\songti,
  box/.style={draw, rounded corners=1.5pt, text width=54mm, minimum height=11mm,
              align=center, line width=0.5pt, inner sep=3pt},
  stage/.style={draw, rounded corners=1.5pt, fill=black!6, text width=14mm,
                minimum height=11mm, align=center, line width=0.5pt, inner sep=2pt},
  ar/.style={->, >=stealth, line width=0.5pt}
]
\node[stage] (s1) {阶段一};
\node[box, right=5mm of s1] (a1) {真实轨迹数据获取与运动学标定};
\node[box, right=6mm of a1] (a2) {DRAC 指标与风险阈值校准};
\draw[ar] (a1) -- (a2);

\node[stage, below=6mm of s1] (s2) {阶段二};
\node[box, right=5mm of s2] (b1) {自研模型与 SUMO 耦合};
\node[box, right=6mm of b1] (b2) {配对种子与多臂对照实验};
\draw[ar] (b1) -- (b2);

\node[stage, below=6mm of s2] (s3) {阶段三};
\node[box, right=5mm of s3] (c1) {可追溯远端消息的控制侧干预};
\node[box, right=6mm of c1] (c2) {风险传播机理与管控启示};
\draw[ar] (c1) -- (c2);

\draw[ar] (a1.south) -- (b1.north);
\draw[ar] (a2.south) -- (b2.north);
\draw[ar] (b1.south) -- (c1.north);
\draw[ar] (b2.south) -- (c2.north);
\end{tikzpicture}

{\zihao{5}\songti 图 1\quad 研究技术路线}
\end{center}

课题推进具备三方面基础条件。数据方面，公开的高分辨率车辆轨迹数据集已完成下载与初步格式解析，可满足单车道标定与多车道扩展的需求。平台方面，自研跟驰模型与 SUMO 的外部步进接口已完成联调，并通过同源数据回放核对，未出现步进丢帧或数值漂移。

前期已完成公开轨迹数据的初步解析、仿真平台联调以及单车道先导实验。在先导条件下，观测到前车紧急制动引发的 DRAC 超阈值事件沿车队向上游逐车传递。该现象表明指标选取与实验通路具备可行性。然而，上述结果局限于单车道、单一源事件与理想通信条件，能否推广至多车道及真实通信环境，尚待后续实验检验。上述前期工作构成后续计划的起点，进度安排中的相应节点指该部分工作的完成与完善时间。
\end{tjformbody}
\tjformsubsection{3．工作进度安排}
\begin{tjformbody}
2026.09.14—2026.10.11：查阅国内外相关文献，了解混合自动驾驶交通、交通风险评价、风险传播及微观交通仿真等方向的研究进展；熟悉已有研究数据、程序与仿真工具，初步明确研究内容和技术路线。

2026.10.12—2026.12.04：开展真实车辆轨迹数据的整理与分析，研究交通风险指标与风险事件识别方法；完成研究方案论证和开题报告，形成初步的数据处理与风险测量方法。

2026.12.05—2027.01.15：构建并完善微观交通仿真环境，设置自动驾驶与人工驾驶车辆的基本行为参数和典型交通场景；结合真实数据对车辆运动状态与风险指标范围进行校核。

2027.01.16—2027.02.28：开展不同混合交通条件及局部扰动条件下的仿真实验，分析风险随时间和空间的变化规律，整理风险传播过程中的主要影响因素和典型现象。

2027.03.01—2027.03.21：根据前期实验结果开展对比实验与敏感性分析，完善风险传播机理分析；整理中期研究成果并完成中期检查材料。

2027.03.22—2027.04.11：探索风险传播抑制与车辆干预控制方法，对比不同控制条件下安全性、效率和运行平稳性等指标，完成主要实验和结果分析。

2027.04.12—2027.04.28：汇总实验结果，完成毕业论文撰写、图表整理和修改，检查程序、实验数据及支撑材料，形成可供检测和评阅的论文定稿。
\end{tjformbody}


\tjformgap{28.68bp}%
\tjformsection[1.5]{三、\tjheadertext 主要参考文献}
\begingroup
\zihao{-4}\songti\selectfont
\baselineskip=23.4bp
\setlength{\bibitemsep}{0pt}%
\setlength{\bibparsep}{0pt}%
\printbibliography[heading=none]
\endgroup

\tjformgap{24.50bp}%
\tjformsection{四、审核意见}
% 本节由导师与专业负责人填写，学生不动。
\tjreviewgap{6.99bp}
\begin{tjreviewbox}[243.6pt]{14.7pt}
  \tjformlabel{指导教师意见：}\par
  \tjformhintbox{（课题难度是否适中、工作量是否饱满、进度安排是否合理、工作条件是否具备、是否同意开题等）}
  \vfill
  \tjformsignature{指导教师签名}{36pt}{295.5pt}
\end{tjreviewbox}

\tjreviewgap{-0.50bp}
\tjreviewtop{6.19bp}
\begin{tjreviewbox}[188.34pt]{12.3pt}
  \tjformlabel{专业审核意见：}\par
  \tjformhintbox[15.6bp]{（适当对各项内容给与意见和评价，字数不限。）}
  \vfill
  \tjformsignature{负责人签名}{42pt}{301.5pt}
\end{tjreviewbox}

\end{document}
